\documentclass[track=junior,level=2,teacher]{kidscode}
\newcommand{\ostar}{\tikz[baseline=-0.6ex]\node[star,star points=5,star point ratio=2.3,draw=black,line width=0.8pt,fill=white,minimum size=0.6cm,inner sep=0pt]{};}
%% a script with a heading: \begin{refscript}{Sprite}{what it does} blocks \end{refscript}
\newenvironment{refscript}[2]{\par\needspace{8\baselineskip}\noindent\textbf{\color{kcaccent}#1:} #2\par\smallskip\begin{center}\begin{scratch}[scale=0.8]}{\end{scratch}\end{center}}
%% the same for a long script: smaller blocks, and the heading stays with the script
\newenvironment{longrefscript}[2]{\par\clearpage\noindent\textbf{\color{kcaccent}#1:} #2\par\smallskip\begin{center}\begin{scratch}[scale=\longscale]}{\end{scratch}\end{center}}
\makeatletter\ifkc@dyslexia\newcommand{\longscale}{0.52}\else\newcommand{\longscale}{0.62}\fi\makeatother
\begin{document}
\kctitle{Final Project Teacher Guide}{My First Arcade Game: milestones, stars, reference game and arcade showcase}
\tableofcontents

%% =====================================================================
\section{The project in one page}
Each child builds a small arcade game in Scratch: a catch game, a maze game, a dodge game or an
idea of their own. They plan it first on a \textbf{game plan card} (who is the player, how do you
win, how do you lose, which sprites, a sketch), build it as Driver and Helper, and then a
\textbf{Player} (a partner, a sibling or you) tests it and finds something to fix. The game must use
five things: \textbf{a forever loop, an if-touching check, a score or lives variable, something
random, and a win or game-over ending}. In week 10 the family plays every game at an arcade
showcase, each child explains one script and one design choice, and receives a certificate.

\begin{center}\small\renewcommand{\arraystretch}{1.35}
\begin{tabularx}{\linewidth}{@{}>{\bfseries\RaggedRight}p{3.0cm}>{\bfseries\RaggedRight}p{3.0cm}X@{}}
\toprule
When & Milestone & What the children do\\
\midrule
Week 9, Warm-up (10 min) & M1 Game plan card & Fill in the game plan card; a partner plays the paper prototype.\\
Week 9, Show me and both Make it blocks (38 min) & M2 Game works & Watch the first must-have being built; build one must-have at a time; at minute 42 a Player playtests; fix one bug from the list; save.\\
Week 9, Share & (check) & Each child shows one must-have that works and says which bug they fixed.\\
Week 10 (the whole session) & M3 Arcade showcase & Last bug check, practice the 1-minute demo, demos and the arcade gallery with sticker feedback, certificates.\\
\bottomrule
\end{tabularx}
\end{center}

\begin{materials}
  \item Print one \textbf{Project Brief} per child (the game plan card page is the one they fill in).
  \item Before week 9: build the reference game yourself (section \enquote{The reference game} below,
        about 30 minutes) and save it as \texttt{Apple-Catch} in the \texttt{Game Makers} folder.
  \item For the core path: an arcade template per child where the player already moves with the
        four arrow keys, saved as \texttt{arcade-} plus the nickname (see below).
  \item Paper, crayons and a few counters or coins for the paper prototype.
  \item Week 10: the week-10 show-me checklist, the sticker feedback cards (last page), stickers,
        one certificate per child (\texttt{extras/\allowbreak{}certificate}), and the progress trackers.
\end{materials}

\begin{teachernote}[Why this project]
The arcade game pulls every Level 2 goal together: loops that keep a game running (J2.G1), if-then
reactions (J2.G2), a score and random surprises (J2.G3), splitting a game into parts (J2.G4, in the
stretch path and in messages), testing and fixing (J2.G5), planning and playtesting with feedback
(J2.G6), and fair, kind games with credit (J2.G7). A game is made for other people to play, which
gives children a real reason to test, listen and improve.
\end{teachernote}

%% =====================================================================
\section{Milestone walkthrough}

\subsection{M1 Game plan card (week 9 Warm-up, about 10 minutes)}
\begin{teacherscript}
\enquote{Today you start your own arcade game. Game makers plan before they build. Our plan is a game plan card, with four questions: who is the player, how do you
win, how do you lose, and which sprites.}
\end{teacherscript}
\begin{enumerate}[itemsep=2pt]
  \item Read the mission and the must-haves aloud together. If you have time (or in the week-8 Share),
        let each child play the reference game for one minute and ask: \enquote{Where was the forever?
        The if touching? The score? What was random? How did it end?}
  \item Show the idea menu. Children choose a catch, maze or dodge game, or their own idea. Offer two
        choices if a child cannot decide.
  \item \textbf{Paper prototype (3--4 minutes):} each child draws the game screen on paper; a partner plays
        it by sliding a coin as the player while the maker moves the \enquote{apple} or \enquote{ghost}
        counter and keeps the score on paper. Ask the Player: \enquote{Did you know how to win?}
  \item Children fill in the game plan card (they can use their rules from week 8): a name for the game, the player and how it moves, how to
        win, how to lose, the sprites they circle, and a sketch. A few words and arrows are perfect.
  \item Point at each must-have icon on the card and ask \enquote{Where is this in your game?} They check the boxes.
\end{enumerate}
\begin{formativecheck}[M1 is done when]
The card has a player, a way to win, a way to lose, the sprites and a sketch, and the child can say
where each must-have will be. Quick question: \enquote{How does someone lose your game?}
\end{formativecheck}

\subsection{M2 Game works (week 9 Show me and Make it, about 38 minutes)}
In Show me, read a sample game plan card aloud (the lesson plan uses Mei's banana game) and build its
first must-have on an arcade template while the children watch.
Driver and Helper swap on the timer as usual. In a pair, they build one game together first (one
child's card), then the other's, or each builds their own game on their own turn. With one child, you
are the Helper and later the Player. Build in this order and play the game after each step:
\begin{enumerate}[itemsep=2pt]
  \item \textbf{The player:} sprite and arrow keys (core path: already there). Play: does it move both ways?
  \item \textbf{The forever loop:} the falling, chasing or moving thing. Play.
  \item \textbf{The if touching:} what happens when they meet. Play.
  \item \textbf{The score (or lives or misses):} \enquote{set score to 0} under the green flag, a
        \enquote{change} inside the if. Play: does it start at 0?
  \item \textbf{Random:} a surprise start place or a surprise wait. Play twice: is it different?
  \item \textbf{The ending:} \enquote{if score = 10 then \ldots{} stop all} and/or the game over. Play to the end.
  \item \textbf{Save:} \menu{File > Save to your computer}, into \texttt{Game Makers}, named with the
        nickname and the game, for example \texttt{Tiger-catch}.
\end{enumerate}
\textbf{Playtest (from minute 42, about 8 minutes).} Pairs swap Macs, or you become the Player.
The Player plays without help while the maker watches with hands in lap and says nothing. The
maker (or you) writes or draws each problem on the playtest bug list in the workbook (or on the
\enquote{My Player found\ldots} lines of the brief).
Then the maker fixes it: change one thing, play again.
\begin{teacherscript}[Ask while you walk around]
\enquote{Which must-have are you building now?} \enquote{Show me the script with the forever.}
\enquote{What will happen when the ghost touches the cat? Let's play and see.} To a Player:
\enquote{Tell the maker what happened, not how to fix it.}
\end{teacherscript}
\begin{formativecheck}[M2 is done when]
All five must-haves work from a fresh green-flag click, a Player has played the game, one thing they
found is fixed, and the game is saved. In Share, each child shows one must-have that works and says
which bug they fixed.
\end{formativecheck}

\subsection{M3 Arcade showcase (week 10, the whole session)}
\begin{enumerate}[itemsep=2pt]
  \item \textbf{Last bug check (Warm-up, 10 minutes):} open the game (\menu{File > Load from your
        computer}), click the green flag twice, play once to win and once to lose. Fix one thing, save, stop.
  \item \textbf{Practice (Show me, 10 minutes):} you model a 1-minute demo, including one design choice;
        each child practices with a partner.
  \item \textbf{Demos (first Make it, 13 minutes):} each child gives the demo and a family member plays one round.
  \item \textbf{Arcade gallery (second Make it, 15 minutes):} visitors play every game and leave stickers;
        in a \enquote{show-me corner} each child shows you a message and a My Block.
  \item \textbf{Celebrate (Share, 10 minutes):} certificates, the stars, showcase safety.
\end{enumerate}
Fill in the week-10 show-me checklist while you watch: most rows show up during the polish, the demo
and the visitors' questions.

%% =====================================================================
\section{Unblocking common problems}
\begin{center}\small\renewcommand{\arraystretch}{1.3}
\begin{longtable}{@{}>{\RaggedRight}p{5.2cm}>{\RaggedRight\arraybackslash}p{\dimexpr\linewidth-5.2cm-2\tabcolsep}@{}}
\toprule
\textbf{Problem} & \textbf{What to do}\\
\midrule\endhead
\bottomrule\endfoot
Cannot choose an idea. & Offer two from the idea menu: \enquote{Catch or dodge?} Then: \enquote{Which sprite do you want to be?}\\
The game plan is huge (ten levels, five enemies). & Circle one player, one thing to catch or dodge, one way to win and one way to lose. The rest goes on a \enquote{later} list.\\
The player does not move. & Click the stage once so Scratch listens to the keys. Check the key in each \enquote{when key pressed} hat. Core path: open the arcade template again.\\
The player or the apple is upside down. & Add \enquote{set rotation style left-right} (blue Motion) under the green flag.\\
The apple falls only once. & The move is under the forever, or there is no forever. Drag \enquote{move 5 steps} inside the forever's mouth.\\
The apple flickers at the top, or misses go up at the start. & The hidden Star marker is touching the top edge. Show the Star, drag it lower and away from the sides, play again.\\
The score races up. & The apple stays touching the bowl. Send it back to the Star straight after \enquote{change score by 1}.\\
The score does not start at 0. & \enquote{set score to 0} must be under the green flag, above the forever.\\
The game ends right at the start. & The variables are reset after the check, or in another sprite. Put every \enquote{set} at the top of the same script, before the forever.\\
The game never ends. & Check the number in \enquote{score = 10} or \enquote{misses = 3} can be reached, and \enquote{stop all} is inside the if.\\
Random does nothing. & The random block is outside the loop, so it picks once. Move it to where the thing restarts.\\
A message does nothing. & Pick the message name from the dropdown in both blocks; check the receiving sprite has the \enquote{when I receive} script.\\
\enquote{It's too hard!} (from the Player) & Change one number: slower falling (move 3), a bigger bowl (costume), more misses before game over. Play again.\\
A library sprite is claimed as their drawing. & Kindly add credit: \enquote{The apple is from the Scratch library.} Painting their own sprite is welcome too.\\
\enquote{I want to add more!} & Great: first all five must-haves, a playtest and a save, then extras one at a time.\\
The project is lost. & \menu{File > Load from your computer}, look in \texttt{Game Makers}, \texttt{Scratch Projects} and Downloads. Save before the wiggle break.\\
Tears or frustration. & Wiggle break. Then play the game together and find one thing that works: \enquote{Look, your bowl already moves both ways!}\\
\end{longtable}
\end{center}

%% =====================================================================
\section{The stars and how to use them kindly}
The brief has four rows, each with one to three stars. The stars are for the child, not a score to
report: there is no total, and children are never compared.

\begin{rubric}[My stars (adult version)]
\small\renewcommand{\arraystretch}{1.35}
\begin{tabularx}{\linewidth}{@{}>{\bfseries\RaggedRight}p{3.4cm}*{3}{>{\RaggedRight\arraybackslash}X}@{}}
 & \ostar & \ostar\,\ostar & \ostar\,\ostar\,\ostar\\
\midrule
It works & The game starts, but parts do not run. & Most of it works; one rule or the ending does not. & Every part works, every time, from a fresh green-flag click.\\ \midrule
It has every must-have & 1 or 2 of the 5. & 3 or 4 of the 5. & All 5: forever, if touching, score or lives, random, a win or game over.\\ \midrule
A player tested it and I fixed something & Someone played the game. & The child watched, listened and named what the Player found. & The child fixed something the Player found and played it again.\\ \midrule
I can say what one script does & Says what the game is about. & Points to a script and names some blocks. & Picks one script and says what each block does, in order.\\
\end{tabularx}
\end{rubric}

\begin{tip}[Using the stars kindly]
\begin{itemize}[itemsep=2pt]
  \item \textbf{The child colors first.} Read each row aloud and let them choose; then say what you saw.
        If you see more stars than they colored, say so. If you see fewer, ask \enquote{Can you show me?}
  \item \textbf{Two moments:} once at the end of week 9 (\enquote{what is left to do?}) and once after
        the showcase (\enquote{look how far you came}).
  \item \textbf{One star is a start.} Pick the one row to work on next and help with it right away.
  \item \textbf{Celebrate the playtest and the fix} as much as the working game: listening to a
        player is a real game maker's skill.
\end{itemize}
\end{tip}

%% =====================================================================
\section{The reference game: \enquote{Apple Catch}}
Build this yourself before week 9 so children can play it, then keep it on \textbf{your} Mac. It is
one complete answer: children make their own game, not a copy. It follows the catch game from weeks 2
and 3 and uses only Level 2 blocks and Grade 2--3 math: directions 0, 90, 180 and 270, whole numbers,
random numbers from 1 to 10, sprites placed by dragging, and \textbf{no negative numbers}. That is why
it counts \textbf{misses up} (game over at 3 misses) instead of taking lives away.

\begin{teachernote}[Why misses and not lives?]
Taking a life away with \enquote{change lives by} needs a negative number, which Grade 2--3 children
have not met. Counting misses up keeps every number positive. A child who prefers lives can use
\enquote{set lives to (lives $-$ 1)} with the green subtraction block and end the game at
\enquote{lives = 0}; both are fine.
\end{teachernote}

\subsection{Setting it up (no blocks yet)}
\begin{enumerate}[itemsep=2pt]
  \item New project (\menu{File > New}). Right-click the Cat and delete it.
  \item Sprites: click \menu{Choose a Sprite} and add the \textbf{Bowl}, the \textbf{Apple} and the
        \textbf{Star}. Drag the Bowl to the bottom middle of the stage. Drag the Star to the
        \textbf{top left}, a little below the top edge and away from the side (the Star marker).
  \item Stage: \menu{Choose a Backdrop}, pick a backdrop you like.
  \item Variables (orange-red \menu{Variables}, \menu{Make a Variable}, \textbf{For all sprites}):
        \texttt{score} and \texttt{misses}. Keep both ticked so they show on the stage.
  \item Save as \texttt{Apple-Catch} in \texttt{Game Makers}.
\end{enumerate}

\subsection{The scripts, block by block}
\begin{refscript}{Star}{The secret start spot. When the flag is clicked, it hides; a hidden sprite is still a place to \enquote{go to}.}
\blockinit{when \greenflag clicked}
\blocklook{hide}
\end{refscript}

\begin{refscript}{Bowl, script 1}{Keeps the bowl the right way up when it turns left.}
\blockinit{when \greenflag clicked}
\blockmove{set rotation style \selectmenu{left-right}}
\end{refscript}

\begin{refscript}{Bowl, scripts 2 and 3}{The player. Each arrow key points the bowl (90 is right, 270 is left), then moves it 10 steps. This is the part the core-path template already has.}
\blockinit{when \selectmenu{right arrow} key pressed}
\blockmove{point in direction \ovalnum{90}}
\blockmove{move \ovalnum{10} steps}
\end{refscript}
\begin{center}
\begin{scratch}[scale=0.8]
\blockinit{when \selectmenu{left arrow} key pressed}
\blockmove{point in direction \ovalnum{270}}
\blockmove{move \ovalnum{10} steps}
\end{scratch}
\end{center}

\begin{refscript}{Apple, My Block \enquote{new apple}}{Sends the apple back to the top with a surprise. It goes to the hidden Star, points right (90) and hops right a random number of times (pick random 1 to 10, each hop 30 steps), then points down (180), ready to fall. This is the random must-have.}
\initmoreblocks{define \namemoreblocks{new apple}}
\blockmove{go to \selectmenu{Star}}
\blockmove{point in direction \ovalnum{90}}
\blockrepeat{repeat \ovaloperator{pick random \ovalnum{1} to \ovalnum{10}}}{
  \blockmove{move \ovalnum{30} steps}
}
\blockmove{point in direction \ovalnum{180}}
\end{refscript}

\begin{refscript}{Apple, the falling script}{Falls forever, as in weeks 2 and 3. Each turn of the loop it moves down 5 steps and asks two questions: caught? (score up, Pop, new apple) and missed? (misses up, new apple).}
\blockinit{when \greenflag clicked}
\blockmove{set rotation style \selectmenu{left-right}}
\blockmoreblocks{new apple}
\blockinfloop{forever}{
  \blockmove{move \ovalnum{5} steps}
  \blockif{if \boolsensing{touching \ovalsensing*{Bowl} ?} then}{
    \blockvariable{change \selectmenu{score} by \ovalnum{1}}
    \blocksound{start sound \selectmenu{Pop}}
    \blockmoreblocks{new apple}
  }
  \blockif{if \boolsensing{touching \ovalsensing*{edge} ?} then}{
    \blockvariable{change \selectmenu{misses} by \ovalnum{1}}
    \blockmoreblocks{new apple}
  }
}
\end{refscript}

\begin{longrefscript}{Bowl, the referee script}{The week-8 referee. It sets the score and misses to 0 \textbf{before} the forever (so the game cannot end at the start), tells the player the goal, then keeps checking: lost? (misses = 3) or won? (score = 10). Each ending says one line and stops everything.}
\blockinit{when \greenflag clicked}
\blockvariable{set \selectmenu{score} to \ovalnum{0}}
\blockvariable{set \selectmenu{misses} to \ovalnum{0}}
\blocklook{say \ovalnum{Catch 10 apples!} for \ovalnum{2} seconds}
\blockinfloop{forever}{
  \blockif{if \booloperator{\ovalvariable{misses} = \ovalnum{3}} then}{
    \blocklook{say \ovalnum{Game over!} for \ovalnum{2} seconds}
    \blockstop{stop \selectmenu{all}}
  }
  \blockif{if \booloperator{\ovalvariable{score} = \ovalnum{10}} then}{
    \blocklook{say \ovalnum{You win!} for \ovalnum{2} seconds}
    \blockstop{stop \selectmenu{all}}
  }
}
\end{longrefscript}

\begin{center}\small\renewcommand{\arraystretch}{1.3}
\begin{tabularx}{\linewidth}{@{}>{\bfseries}p{4.6cm}X@{}}
\toprule
Must-have & Where it is in the reference game\\
\midrule
A forever loop & The Apple's falling script (and the Bowl's referee keeps checking forever).\\
An if-touching check & \enquote{touching Bowl?} (a catch) and \enquote{touching edge?} (a miss).\\
A score or lives variable & \texttt{score} goes up by 1 for each catch; \texttt{misses} goes up by 1 for each miss.\\
Random & \enquote{pick random 1 to 10} hops in the \enquote{new apple} block: the apple falls from a different place each time.\\
A win or game-over ending & The Bowl's referee: \enquote{score = 10} says You win!, \enquote{misses = 3} says Game over! Both end with \enquote{stop all}.\\
Stretch: a My Block & \enquote{new apple}, used three times.\\
\bottomrule
\end{tabularx}
\end{center}

\begin{tip}[Questions to ask about the reference game]
\enquote{How do we make the game harder?} (Move 7 steps instead of 5, or win at 15.)
\enquote{Why are the \enquote{set} blocks above the referee's forever?} (So the game starts at 0 before it checks for a win: the week-8 bug.)
\enquote{What would happen without \enquote{new apple} after a catch?} (The score races up: the week-3 bug.)
\enquote{Where is the surprise?} (The random hops.)
\enquote{Why count misses instead of taking lives away?} (No numbers below zero.)
\enquote{Whose pictures and sounds are these?} (The Scratch library: we say so at the showcase.)
\end{tip}

\subsection{Without the My Block (core version)}
My Blocks are week 7 and the stretch path. A child who does not use one puts the same four blocks
(\enquote{go to Star}, \enquote{point in direction 90}, the random repeat, \enquote{point in direction 180})
in each place where the reference game says \enquote{new apple}: above the forever and inside both ifs.
Another random surprise that needs no copies: \enquote{wait (pick random 1 to 10) seconds} after \enquote{go to Star},
so the apple waits a surprise time at the top (Mei's idea in the week-9 lesson).
Right-click \enquote{go to Star} and choose \menu{Duplicate} to copy the four blocks at once. This is a
good moment to show the stretch: \enquote{Three copies of the same four blocks: what if we gave them a name?}

\subsection{The core-path arcade template}
Make it before week 9, about two minutes per child: a new project with a backdrop and one player sprite
(the Bowl for a catch game, the Cat for a maze or dodge game) that already moves with the four arrow
keys. Each key has its own script: \enquote{point in direction} 0 (up arrow), 90 (right arrow), 180
(down arrow) or 270 (left arrow), then \enquote{move 10 steps}; a green-flag script sets the rotation
style to left-right. Save one copy per child in \texttt{Game Makers} as \texttt{arcade-} plus the
nickname, for example \texttt{arcade-Tiger}. In week 9 the child opens their copy and adds one
must-have at a time; for a catch game they add the Star marker first.

\subsection{The stretch path: a My Block or a second level}
\textbf{A My Block:} the \enquote{new apple} block above, or \enquote{celebrate} (say, sound, next
costume) used for every catch. Check it is used in two or more places. \textbf{A second level:} inside
the \enquote{score = 10} if, instead of stopping, \enquote{broadcast level 2}; the Stage receives it and
switches the backdrop, and the Apple receives it, sets the score back to 0 and moves faster from then on
(for example a \texttt{speed} variable used in \enquote{move speed steps}, set to 5 at the start and 8
for level 2). Keep \enquote{stop all} for the win at the end of level 2.

%% =====================================================================
\section{Showcase plan (week 10)}

\subsection{The arcade gallery}
The showcase is at home, for family only. Each game gets an \enquote{arcade machine}: a Mac with the
game open and the stage in full screen (the button above the stage), the game plan card beside it,
and a folded card with the game's name, the maker's \textbf{nickname} and \enquote{How to play} in one
line (for example \enquote{Arrow keys. Catch 10 apples. 3 misses and it is over.}). Visitors move from
machine to machine. Twenty to twenty-five minutes is plenty with four children.

\begin{center}\small\renewcommand{\arraystretch}{1.3}
\begin{tabularx}{\linewidth}{@{}>{\bfseries}p{3.0cm}X@{}}
\toprule
Role & Job\\
\midrule
Game maker & Gives the 1-minute demo, clicks the flag for each visitor, answers one question.\\
Helper & (A partner or a sibling.) Hands out feedback cards and stickers, reads the \enquote{How to play} card aloud.\\
Visitors & Read the rules, play one game, ask one question (\enquote{Why did you choose\ldots?}), put stickers on a feedback card.\\
Host (you) & Welcomes the family, keeps time, makes sure every machine gets visitors, reads the safety reminder, plays every game once.\\
\bottomrule
\end{tabularx}
\end{center}

\subsection{The 1-minute demo script}
Write it on a sticky note for each child. Practice once before the family arrives.
\begin{enumerate}[itemsep=2pt]
  \item \enquote{Hi! My game is called \underline{\hspace{3cm}}.}
  \item \enquote{You are the \underline{\hspace{2.5cm}}. You win when \underline{\hspace{3cm}}. You lose when \underline{\hspace{3cm}}.}
  \item \enquote{Watch: I click the green flag.} (Play for about 20 seconds.)
  \item Point to one script: \enquote{This script makes \underline{\hspace{3cm}}. It goes forever: first \underline{\hspace{2.5cm}}, then if \underline{\hspace{2.5cm}}.}
  \item One design choice: \enquote{I chose \underline{\hspace{3cm}} because \underline{\hspace{3.5cm}}.}
  \item \enquote{My Player found \underline{\hspace{3cm}} and I fixed it.}
  \item \enquote{The sprites and sounds are from the Scratch library. Do you want to play?}
\end{enumerate}

\subsection{Certificates and closing}
Read the feedback cards aloud to each child, then hand out the certificate (write the nickname or
first name by hand). Close with a round of \enquote{one thing I am proud of, one game I want to make
next}, and color the last stars on the progress tracker. Remind everyone: the arcade closes when the
timer rings. Good games end.

\begin{safety}[Showcase safety]
\begin{itemize}[itemsep=1pt]
  \item Family only, at home, on the screen. Nothing is uploaded, posted or emailed.
  \item Nicknames or first names only on the screen and on the arcade cards.
  \item Photos only if the child agrees, with no names on the screen, and kept in the family.
  \item Credit: every game says the sprites and sounds come from the Scratch library.
  \item Ask first: a child decides who may play their game, and nobody changes someone else's game.
\end{itemize}
\end{safety}

\begin{teachernote}[Fitting it into the 60-minute session]
0--10 last bug check and save (warm-up); 10--20 model and practice the 1-minute demo (show me);
20--33 demos while the family plays, week-10 checklist (make it); 33--35 wiggle break (act out every
game); 35--50 arcade gallery with sticker feedback and the show-me corner (make it); 50--60
certificates, stars, showcase safety (share).
\end{teachernote}

\clearpage
\section{Sticker feedback cards (print and cut)}
Visitors put a sticker next to what they liked, and a grown-up can write one more idea.
\newcommand{\feedbackcard}{%
%% alt: A feedback card that says Your game! with five dashed sticker spots next to fun to play, the rules were clear, it was fair, the sounds and the surprise, and a line for one more idea.
  \begin{tikzpicture}[spot/.style={circle,draw=kcgray,dashed,minimum size=0.8cm,inner sep=0pt}]
    \draw[kcpurple,line width=1.2pt,rounded corners=6pt] (0,0) rectangle (8.2,8.0);
    \node[anchor=north west,font=\bfseries] at (0.3,7.8) {\faGamepad\ Your game!};
    \node[anchor=north west,font=\small] at (0.3,7.1) {Put a sticker by what you liked.};
    \foreach \icon/\word [count=\i] in {\faSmile/fun to play,\faListOl/the rules were clear,\faBalanceScale/it was fair,\faVolumeUp/the sounds,\faDice/the surprise}
      {\node[spot] at (0.9,6.7-\i*1.0) {};
       \node[anchor=west] at (1.6,6.7-\i*1.0) {\icon\ \word};}
    \node[anchor=north west,font=\small] at (0.3,1.05) {One more idea:};
    \draw[kcgray!60] (0.3,0.3) -- (7.9,0.3);
  \end{tikzpicture}}
\begin{center}
\feedbackcard\hfill\feedbackcard\par\bigskip
\feedbackcard\hfill\feedbackcard
\end{center}
\end{document}
